\chapter{Template Guide}%
\label{s_template_guide}

This section explains how to use the thesis template.

\section{Template Setup}%
\label{s_setup}

The template is loaded by specifying \lstinline|\documentclass{fsdthesis}| at
the start of the code.  Note that options provided in square brackets are
global, they will be passed to all loaded packages.
Thus, the initialization of the \lstinline|fsdthesis| class is deferred.
Options can be specified using the \lstinline|\fsdthesisSetup| command as shown
in \cref{l_template_setup}:
\begin{itemize}
  \item \lstinline|language| determines the language of labels and standard
        headings throughout the document and sets up the hyphenation rules.  At
        the moment it does not trigger automatic translation; this feature may
        be considered in future versions.  The \lstinline|abstract|
        environment, which is always included in German and English versions,
        locally overrides this option.
        At the time of writing, the option only accepts the values
        \lstinline|en| for English and \lstinline|de| for German.  This
        limitation is acceptable since it corresponds to \gls{TUM} regulations.
  \item \lstinline|type| may be \lstinline|MSc| (Master's Thesis),
        \lstinline|BSc| (Bachelor's Thesis), \lstinline|semester| (semester
        thesis) or \lstinline|IDP| (interdisciplinary project).  Further types
        may be added in the future.
  \item \lstinline|title| specifies the main title displayed on the cover page
        and in the footers and recorded in the \gls{PDF} metadata.  This title should
        be given in the language of the main text.
  \item \lstinline|subtitle| specifies a subtitle, which is commonly a
        translated version of the main title.
  \item \lstinline|authorLabelForm| provides a choice between \lstinline|f|
        (female), \lstinline|m| (male) and \lstinline|d| (diverse) forms of the
        \emph{author} label on the cover page.
  \item \lstinline|author| sets the author name to be displayed on the cover
        page, in the footer, in the statutory declaration and in the \gls{PDF}
        metadata.
  \item \lstinline|authorID| gives the student number of the author.
  \item \lstinline|supervisorsLabelForm| works like the
        \lstinline|authorLabelForm|, and supports additional uppercase values
        that represent plural forms.
  \item \lstinline|supervisors| lists the supervisors, separated by
        \lstinline|\\|; these are printed on the cover page.
  \item \lstinline|examinersLabelForm| and \lstinline|examiners| allow
        overriding the predefined examiner, e.g., in case of additional
        examiners.  Usually, these options should not be specified.
  \item \lstinline|date| should correspond to the submission date and is stated
        on the cover page and in the statutory declaration.
  \item \lstinline|anonymize| allows to generate an anonymized version of the
        thesis for automated plagiarism checks.  This erases the personal
        information of author, supervisors and examiner from the document.
        However, if names are manually entered in the text, these are not
        removed.
\end{itemize}
These options are processed at \lstinline|\begin{document}|, i.e., after custom
code in the preamble.

\lstinputlisting[
  language={[LaTeX]TeX},
  linerange=fsdthesisSetup-fsdthesisSetup,
  rangebeginprefix=\%<,
  rangebeginsuffix=>,
  rangeendprefix=\%</,
  rangeendsuffix=>,
  includerangemarker=false,
  float,
  caption={Template setup code (\texttt{Thesis.tex})},
  label={l_template_setup}
]{Thesis.tex}

\section{Content}

\subsection{Front Matter}

The front matter is delimited by \lstinline|\begin{frontMatter}| and
\lstinline|\end{frontMatter}|.  It starts with the statutory declaration and
abstracts in English and German.  These are followed by the table of contents
(\lstinline|\makeTableOfContents|), the list of figures
(\lstinline|\makeListOfFigures|), the list of acronyms
(\lstinline|\makeListOfAcronyms|) and the table of symbols defined in
\verb|TableOfSymbols.tex|.

\subsection{Main Content}

The structure of the main content depends on the specific thesis project.
Please refer to the thesis guide~\cite{FSD_Thesis_Guide} for a typical outline
and talk to your supervisors.

Generally, it is a good idea to split at least the main chapters into a few
files and include these with the \lstinline|\input{}| command.  The document
structure is determined by sectioning commands such as \lstinline|\chapter{}|,
\lstinline|\section{}|, \lstinline|\subsection{}| and
\lstinline|\subsubsection{}|.
Ideally, each of these is followed by a \lstinline|\label{}| to enable
cross-references, see also \cref{s_cref}.  The main content is placed between
\lstinline|\begin{mainContent}| and \lstinline|\end{mainContent}|.

\subsection{Back Matter}

Between \lstinline|\begin{backMatter}| and \lstinline|\end{backMatter}|,
optional acknowledgements may be included (see \verb|Acknowledgements.tex|),
followed by the bibliography (\lstinline|\makeBibliography|) and optional
appendices.  The latter are declared as chapters in \verb|Appendix.tex|.

\section{Special Effects}

\subsection{Acronyms}

On first use of an acronym, the term needs to be spelled out.  To automate this
behavior, define acronyms in \verb|Acronyms.tex|.

Acronyms can be used as shown in the following example, where the definition
has been included for educational purposes:
\begin{quote}
%<acronymExample>
\newacronym{LQR}{LQR}{Linear-Quadratic Regulator}
The \gls{LQR} minimizes a quadratic performance index and results
in a linear state feedback law.  The \gls{LQR} is thus an optimal
controller.  \Glspl{LQR} provide an intuitive way to specify
desired dynamics.
%</acronymExample>
\end{quote}

The code for this example is given in \cref{l_acronym_example}.  The first
argument to \lstinline|\newacronym| defines an identifier, which can deviate
from the short form given in the second argument.  The \lstinline|\gls| command
works with this identifier.
At the start of a sentence, \lstinline|\Gls| should be used to capitalize the
first letter of the acronym.  Plural forms can be obtained using
\lstinline|\glspl| and \lstinline|\Glspl|.

\FloatBarrier{}
\lstinputlisting[
  language={[LaTeX]TeX},
  linerange=acronymExample-acronymExample,
  rangebeginprefix=\%<,
  rangebeginsuffix=>,
  rangeendprefix=\%</,
  rangeendsuffix=>,
  includerangemarker=false,
  float,
  caption={Acronym example (\texttt{\currfilename})},
  label={l_acronym_example}
]{\currfilename}

\subsection{Math}

The following examples are generated by \cref{l_math_example}.
Individual equations can be referenced if the are labeled appropriately, such
as \cref{eq_T_f,eq_x_u,eq_demoivre}.

{ % local scope for new macros
%<mathExample>
\newcommand{\Vector}[1]{\bm{#1}}
\newcommand{\Vx}{\Vector{x}}
\newcommand{\Vy}{\Vector{y}}
\newcommand{\Vu}{\Vector{u}}
\newcommand{\Vf}{\Vector{f}}
\newcommand{\Real}{\mathbb{R}}

A nonlinear autonomous dynamic system may be defined by its dynamics
$\dot{\Vx}(t)=\Vf(\Vx(t), \Vu(t))$, where:
\begin{subequations}
\begin{align}
  T &= [0, t_f]
  &
  \Vf&: \Real^n \times \Real^m \rightarrow \Real^n,~
       (\Vx, \Vu) \mapsto \Vf(\Vx, \Vu)
  \label{eq_T_f}
  \\
  \Vx&: T \rightarrow \Real^n,~ t \mapsto \Vx(t)
  &
  \Vu&: T \rightarrow \Real^n,~ t \mapsto \Vu(t)
  \label{eq_x_u}
\end{align}
\end{subequations}

A second-order lag can be written in the state-space form $\dot{\Vx} = \bm{A}
\Vx + \bm{B} \Vu$ with the following constants:
\begin{align}
  \bm{A} &= \begin{bmatrix}
              0 & 1 \\ -\omega_0^2 & -2\,\zeta\,\omega_0
            \end{bmatrix}
  &
  \bm{B} &= \begin{bmatrix} 0 \\ k \,\omega_0^2 \end{bmatrix}
\end{align}

The exponential of an imaginary number can be expressed by
\begin{equation}
  \label{eq_demoivre}
  \exp(\iu \varphi) = \cos(\varphi) + \iu \sin(\varphi)
  .
\end{equation}
%</mathExample>
}

\lstinputlisting[
  language={[LaTeX]TeX},
  linerange=mathExample-mathExample,
  rangebeginprefix=\%<,
  rangebeginsuffix=>,
  rangeendprefix=\%</,
  rangeendsuffix=>,
  includerangemarker=false,
  float,
  caption={Math example (\texttt{\currfilename})},
  label={l_math_example}
]{\currfilename}

Common special symbols are produced according to \cref{l_math_symbols}:
%<mathSymbols>
\begin{align}
  \sum_{k=-\infty}^{\infty},
  \sum\nolimits_{k=-\infty}^{\infty},
  \prod_{k=-\infty}^{\infty},
  \prod\nolimits_{k=-\infty}^{\infty},
  \int\limits_{t=-\infty}^{\infty},
  \int_{t=-\infty}^{\infty},
  \iint,
  \iiint,
  \oint,
  \nonumber\\
  \sin, \cos, \tan, \arcsin, \arccos, \arctan, \nonumber\\
  \sinh, \cosh, \tanh, \arcsinh, \arccosh, \arctanh, \nonumber\\
  \exp, \log, \min, \max, \arg, \lim_{x\to 0} \nonumber\\
  \begin{bmatrix}
    a_{11} & \cdots & a_{1n} \\
    \vdots & \ddots & \vdots \\
    a_{m1} & \cdots & a_{mn}
  \end{bmatrix},
  \begin{pmatrix} v_1 \\ \vdots \\ v_m \end{pmatrix},
  |x|, \Vert x\Vert^2_2, \left\Vert\frac{x}{y}\right\Vert^2_2, \sqrt[3]{x},
  \nabla^2_{xy}f = \frac{\partial^2 f}{\partial y \partial x} =
  \mathcal{D}^2_{xy}f
  \nonumber\\
  <, \le, =, \ge, >,
  \approx, \overset{!}{=}, \eqqcolon, \coloneqq,
  \leftarrow, \leftrightarrow, \rightarrow,
  \Leftarrow, \Leftrightarrow, \Rightarrow,
  \subsetneq, \subset, \subseteq, \supseteq, \supset, \supsetneq,
  \cap, \cup, \setminus,
  \in, \exists, \forall \nonumber\\
  \mathcal{L}, \mathcal{F}, \mathfrak{R}, \mathfrak{I},
  \mathbb{N}, \mathbb{Q}, \mathbb{R}, \mathbb{C} \nonumber\\
  \bm{x}^\T, \bm{x}^\HT \nonumber
\end{align}
%</mathSymbols>

\lstinputlisting[
  language={[LaTeX]TeX},
  linerange=mathSymbols-mathSymbols,
  rangebeginprefix=\%<,
  rangebeginsuffix=>,
  rangeendprefix=\%</,
  rangeendsuffix=>,
  includerangemarker=false,
  float,
  caption={Math symbols (\texttt{\currfilename})},
  label={l_math_symbols}
]{\currfilename}

Spaces of varying width are produced by \lstinline|\,|, \lstinline|\;|,
\lstinline|~|, \lstinline|\quad|, \lstinline|\qquad|.

\subsection{Units}

%<unitsExample>
The use of units is demonstrated by \cref{l_units_example}.
For example, the standard air density at mean sea level is
\SI{1.225}{kg.m^{-3}}.
This also works in math mode: $\rho_0 = \SI{1.225}{kg.m^{-3}}$.
%</unitsExample>

\lstinputlisting[
  language={[LaTeX]TeX},
  linerange=unitsExample-unitsExample,
  rangebeginprefix=\%<,
  rangebeginsuffix=>,
  rangeendprefix=\%</,
  rangeendsuffix=>,
  includerangemarker=false,
  float,
  caption={Units example (\texttt{\currfilename})},
  label={l_units_example}
]{\currfilename}

\subsection{Citations}

%<citationExample>
Citations are an integral part of good academic
practice~\cite[\S\,11\,(5)\,4., p.\,5]{Good_Academic_Practice}:
\begin{quote}
  \textsuperscript{1}%
  Findings and ideas of other researchers as well as relevant
  publications of other authors should be cited
  (see TUM Citation Guidelines). [\ldots]
\end{quote}
The \gls{TUM} statutes also require that the
\textquote{source of any data, organisms, materials or software
used in the research process is designated as such}~%
\cite[\S\,7\,(2), p.\,3]{Good_Academic_Practice}.
%</citationExample>
This applies not only to direct citations as shown above and in the associated
\cref{l_citation_example}.
Indirect quotations, i.e., paraphrased ideas, also need to be cited
appropriately!

To conveniently handle your sources, choose an appropriate software package for
literature management and export its database to the \texttt{References.bib}
file associated with your document.
Each entry is identified by a unique name, such as \texttt{ENTRY}.
In the \LaTeX{} source code, use \lstinline|\cite[pp.\,X--Y]{ENTRY}|,
\lstinline|\citeauthor*{ENTRY}|, or \lstinline|\citetitle{ENTRY}| (in rare
cases) to create a reference to the entry and list it in the bibliography (see
p.\,\pageref{s_references}).
Unused entries are not listed.

\lstinputlisting[
  language={[LaTeX]TeX},
  linerange=citationExample-citationExample,
  rangebeginprefix=\%<,
  rangebeginsuffix=>,
  rangeendprefix=\%</,
  rangeendsuffix=>,
  includerangemarker=false,
  float,
  caption={Citation example (\texttt{\currfilename})},
  label={l_citation_example}
]{\currfilename}

\subsection{Figures}

In \LaTeX{}, the \lstinline|figure| environment is a so-called \emph{float},
i.e., a block that breaks the text flow.
In general you should let \LaTeX{} decide where exactly to place figures,
instead of enforcing a specific location; if needed, provide hints on the
desired placement\footnote{%
\url{https://www.latex-project.org/publications/2014-FMi-TUB-tb111mitt-float-placement.pdf}}.

Use vector graphics if possible.
Commonly-used file formats include \gls{EPS} and \gls{PDF}.
As an example, see \cref{f_figure_example} and the associated
\cref{l_figure_example}.

%<figureExample>
\begin{figure}
  \centering
  \includegraphics[height=30mm]{graphics/Logo-FSD-black.pdf}
  \caption{FSD Logo}%
  \label{f_figure_example}
\end{figure}
%</figureExample>

\lstinputlisting[
  language={[LaTeX]TeX},
  linerange=figureExample-figureExample,
  rangebeginprefix=\%<,
  rangebeginsuffix=>,
  rangeendprefix=\%</,
  rangeendsuffix=>,
  includerangemarker=false,
  float,
  caption={Figure example (\texttt{\currfilename})},
  label={l_figure_example}
]{\currfilename}

\subsection{Tables}

Tables, like figures, are \emph{floating} environments.
Let \LaTeX{} decide their placement unless the results turn out problematic.
Rows and columns are provided by the \lstinline|tabular| environment, which is
usually used inside a \lstinline|table|.
For example, \cref{l_table_example} creates \cref{t_table_example}.

%<tableExample>
\begin{table}
  \caption{Exemplary table/tabular environment}%
  \label{t_table_example}
  \centering
  \begin{tabular}{ccc}
    \toprule
    A & B & C \\
    \midrule
    a & b & c \\
    d & e & f \\
    \bottomrule
  \end{tabular}
\end{table}
%</tableExample>

\lstinputlisting[
  language={[LaTeX]TeX},
  linerange=tableExample-tableExample,
  rangebeginprefix=\%<,
  rangebeginsuffix=>,
  rangeendprefix=\%</,
  rangeendsuffix=>,
  includerangemarker=false,
  float,
  caption={Table/tabular example (\texttt{\currfilename})},
  label={l_table_example}
]{\currfilename}

\subsection{Algorithms}

Mathematical algorithms can be presented in the form of pseudo-code.
For example, \cref{l_algorithm_example} creates \cref{a_algorithm_example}.

%<algorithmExample>
\begin{algorithm}
  \caption{Exemplary algorithm}%
  \label{a_algorithm_example}
  \begin{algorithmic}
    \Procedure{ExplicitRungeKutta}{%
        $x_0$, $f$, $h$, $a_{i,j}$, $b_j$, $q$, $N$}
      \For{$k \leftarrow 1,\ldots,N$}
        \For{$i \leftarrow 1,\ldots,q$}
          \State{$s \leftarrow x_{k-1} +
                  h \sum\limits_{j=1}^{i-1} r_j a_{i,j}$}
          \State{$r_i \leftarrow f(s)$}
        \EndFor{}
        \State{$x_k \leftarrow x_{k-1} +
                  h \sum\limits_{j=1}^{q} r_j b_j$}
      \EndFor{}
      \State{\Return{$(x_0, \ldots, x_N)$}}
    \EndProcedure{}
  \end{algorithmic}
\end{algorithm}
%</algorithmExample>

\lstinputlisting[
  language={[LaTeX]TeX},
  linerange=algorithmExample-algorithmExample,
  rangebeginprefix=\%<,
  rangebeginsuffix=>,
  rangeendprefix=\%</,
  rangeendsuffix=>,
  includerangemarker=false,
  float,
  caption={Algorithm example (\texttt{\currfilename})},
  label={l_algorithm_example}
]{\currfilename}


\subsection{Listings}

\Cref{l_hello_world} includes a hello-world program from an external file.
The \LaTeX{} code for this example is given by \cref{l_listing_example}.
You can find more elaborate examples in the source code of this document
(\texttt{\currfilename}).

%<listingExample>
\lstinputlisting[
  language=matlab,
  float,
  caption={%
    Hello world program (\texttt{listings/hello\_world.m})
  },
  label={l_hello_world}
]{listings/hello_world.m}
%</listingExample>

\lstinputlisting[
  language={[LaTeX]TeX},
  linerange=listingExample-listingExample,
  rangebeginprefix=\%<,
  rangebeginsuffix=>,
  rangeendprefix=\%</,
  rangeendsuffix=>,
  includerangemarker=false,
  float,
  caption={Listing example (\texttt{\currfilename})},
  label={l_listing_example}
]{\currfilename}

\subsection{Cross-References}%
\label{s_cref}

Mark the target with \lstinline|\label{id}|, where \lstinline|id| is a custom
identifier.  Place the label directly after the relevant sectioning command or
the caption of a table or figure, or within an equation.  Then use
\lstinline|\cref{id}| to print a reference that states the type of the target
and its number.  At the start of a sentence, use \lstinline|\Cref{id}| instead.
To print only the number of the target, use the low-level \lstinline|\ref{id}|
command.

For example, consider the following equation:
%<crefExampleEq>
\begin{align}
  \label{eq_lift}
  L &= \frac{1}{2}\, \rho\, V_A^2\, S\, C_L
\end{align}
\Cref{eq_lift} expresses the relation between the lift force $L$
and the lift coefficient $C_L$ with the density $\rho$,
\gls{TAS} $V_A$ and reference area $S$.
%</crefExampleEq>

The code for this example is given in \cref{l_cref_example_eq}.

\lstinputlisting[
  language={[LaTeX]TeX},
  linerange=crefExampleEq-crefExampleEq,
  rangebeginprefix=\%<,
  rangebeginsuffix=>,
  rangeendprefix=\%</,
  rangeendsuffix=>,
  includerangemarker=false,
  float,
  caption={Cross-reference example},
  label={l_cref_example_eq}
]{\currfilename}

\section{Notes}

You may use the commands \lstinline|\todo{note 1}|\todo{note 1} and
\lstinline|\todo[inline]{note 2}|\todo[inline]{note 2} to place notes in the
margin and in the main text.

\section{Compilation}

This template is designed and tested with the \lstinline|xelatex| compiler.
To make acronyms and citations work correctly, call the compiler via the
\lstinline|latexmk| tool; an appropriate configuration file is provided with
the template.
On a unixoid system you may simply call \lstinline|make|.

If you are working with the \emph{TeXstudio} application, you may load the
configuration defined in \texttt{Thesis.txsprofile} via \emph{Load profile} in
the \emph{Options} menu.
This sets up the appropriate commands.

Working with \emph{MiKTeX} on \emph{Windows}, the \lstinline|latexmk| tool may
fail to run, lacking a \lstinline|perl| script
interpreter\footnote{\url{https://www.perl.org/}}.
This can be solved by either installing a suitable \emph{Perl} distribution or
simply using the \emph{Perl} environment shipped with \emph{Matlab}, which is
most probably installed on your computer anyway.
In this case, edit the environment variable \lstinline|PATH| and add an entry
pointing to your \emph{Matlab} installation, such as
\lstinline|C:\Program Files\MATLAB\R2023a\sys\perl\win32\bin\| --
replacing \lstinline|R2023a| by your preferred \emph{Matlab} release, of
course.
